\title{Parameter-Efficient Orthogonal Finetuning via Factorization}

\begin{abstract}
As large foundation models become increasingly prevalent, the cost associated with training them from the ground up grows ever more prohibitive. Consequently, finding effective strategies for customizing these models for specific downstream tasks is vital. This work explores Orthogonal Finetuning (OFT), a theoretically-grounded strategy for transferring foundation models to new tasks. Although OFT boasts impressive generalization, it remains parameter-heavy, largely due to the intrinsic high dimensionality of orthogonal matrices. To tackle this issue, we analyze OFT through the lens of information transmission, extracting core principles conducive to heightened parameter efficiency. Motivated by the computational benefits of the Cooley-Tukey fast Fourier transform algorithm in efficient information processing, we introduce a streamlined orthogonal parameterization based on butterfly structures. This new approach, when integrated with OFT, yields Orthogonal Butterfly (BOFT), a novel, parameter-efficient finetuning method. Not only does BOFT encompass OFT as a particular case, but it also presents a more flexible, general orthogonal finetuning paradigm. We validate our method through thorough experiments across a range of vision transformers, large language models, and text-to-image diffusion models, evaluating adaptation to diverse tasks in both vision and language domains.
\end{abstract}

\section{Introduction}

Cutting-edge systems like ChatGPT~\cite{brown2020language,chen2021evaluating} and Stable Diffusion~\cite{rombach2022high} illustrate the extraordinary generalization potential achieved by large-scale foundation models. The escalation in these models' capabilities, however, is accompanied by a dramatic surge in their parameter counts

(\emph{e.g.}, GPT-3 comprises roughly 175B parameters), presenting significant obstacles to training such systems from scratch. As a result, the broader advancement of the field increasingly depends on the development of approaches for repurposing existing foundation models for new tasks, without resorting to complete retraining. Achieving this calls for efficient techniques to tailor pre-trained models to downstream scenarios. Task adaptation is mainly accomplished through three strategies: \emph{model finetuning}~\cite{hulora2022,zaken2022bitfit,guo2020parameter,raffel2020exploring,qiu2023controlling,zhang2022adaptive,chavan2023one}, which retains the original model architecture and updates a subset of parameters; \emph{adapter tuning}~\cite{rebuffi2017learning,houlsby2019parameter,pfeiffer2020adapterfusion,lin2020exploring,he2021towards}, which augments the model with new trainable modules; and \emph{prompt tuning}~\cite{li2021prefix,lester2021power}, which introduces additional tunable prefix tokens to the input stream. Of these, model finetuning is especially notable for its simplicity, practical efficacy, and for leaving inference speed unaltered.

Model finetuning fundamentally relies on maintaining similarity between the pretrained and finetuned models, as measured by specific criteria, in order to retain the knowledge acquired during pretraining. Common practice in existing finetuning approaches is to use a low learning rate, an intuitive strategy that limits the deviation between the two models. Given the inherent structure present in weight matrices, more methodical techniques focus on preserving the relational structure—specifically, the pairwise angles among neurons—within these matrices. Building on this, Orthogonal Finetuning~(OFT)~\cite{qiu2023controlling} has emerged as an innovative finetuning framework wherein all neurons in a given layer are subjected to the same orthogonal transformation. This mechanism guarantees that the pairwise angles between neurons are exactly maintained throughout finetuning. OFT has shown strong results in terms of both generalization and convergence when applied to finetuning text-to-image diffusion models~\cite{qiu2023controlling}. Nevertheless, due to the typically large dimensionality of orthogonal matrices, OFT can require a substantial number of trainable parameters. To improve parameter efficiency, OFT employs a block-diagonal orthogonal structure, thereby decreasing parameter count. Yet, this increase in efficiency incurs certain costs—the orthogonal transformations are now limited to operate within fixed blocks, as dictated by the predetermined block-diagonal sparsity pattern. While this approach does reduce parameter requirements, it can induce certain undesirable inductive biases (\emph{e.g.}, limiting expressivity and rendering the matrix unable to approximate standard linear transformations). Furthermore, determining the optimal sparsity pattern for such block structures remains an unresolved challenge.

A central challenge lies in constructing a dense orthogonal matrix in a way that is efficient with respect to the number of parameters. Typically, a fully dense $d$-dimensional orthogonal matrix requires on the order of $\mathcal{O}(d^2)$ parameters, which appears computationally expensive. However, we propose a novel strategy: rather than parameterizing the entire matrix directly, we construct the dense orthogonal matrix through a product of multiple sparse, factorized matrices. This perspective leverages the concept that parameter count can be minimized by increasing the computational steps required—specifically, by multiplying the sparse factors at each iteration during training. 

Viewing this process as akin to an information transmission problem allows us to frame the composition of a dense orthogonal matrix as propagating signals through a grid-structured network. In this context, building the matrix via successive multiplication of sparse factors is analogous to transmitting information across nodes of the grid. This formulation provides a flexible framework for designing sparse matrix factorization schemes that control parameter counts while preserving sufficient expressivity to represent dense orthogonal transformations.

For further parameter reduction, our framework draws conceptual parallels with the butterfly networks used in the Cooley-Tukey fast Fourier transform algorithm~\cite{cooley1965algorithm}, renowned for facilitating efficient information mixing between nodes~\cite{klasing1994broadcasting}. The butterfly graph induces a specific structure of sparse matrix products—known as butterfly factorization. Utilizing this approach, one can express a $d$-dimensional dense matrix as a product of $\mathcal{O}(\log d)$ sparse matrices, each containing $\mathcal{O}(d)$ nonzero elements. When these sparse factors are all orthogonal, the total number of parameters required drops dramatically to $\mathcal{O}(d\log d)$—a substantial decrease from the quadratic scaling of naive parameterization.

Building upon this efficient parameterization, we introduce Orthogonal Butterfly~(BOFT), a novel finetuning method that is both parameter-efficient and flexible. BOFT encompasses OFT as a special instance, extending to a broader family of orthogonal finetuning methods. Notably, both the block-diagonal and butterfly architectures share the key feature of imposing sparsity, thereby reducing the trainable parameter space. Comparing our method to LoRA~\cite{hulora2022}, which utilizes low-rank decomposition, highlights an important distinction: both low-rank and sparse representations are fundamental structured matrix families~\cite{candes2011robust} employed to realize parameter-efficient adaptations.

In contrast to OFT, which employs a block-diagonal structure to balance between model expressivity and parameter regularity, BOFT leverages the butterfly structure to achieve a more continuous transition between matrices drawn from the full orthogonal group (which supports expressivity) and the identity matrix (which supports regularity). This approach allows us to define a more compact hypothesis space within the orthogonal group, improving suitability for various downstream applications. Notably, the butterfly structure underpins many widely used fast linear transforms, including the discrete Fourier and discrete cosine transforms. We therefore posit that employing this structure imparts a natural inductive bias, fostering better generalization and reducing the risk of overfitting. Our rationale stems from the fact that butterfly matrices can efficiently represent many classical linear transformations, a capability the block-diagonal construction of OFT fundamentally lacks. Our key contributions are as follows:

\begin{itemize}[leftmargin=*,nosep]
    \item We introduce a new perspective on parameter efficiency in orthogonal finetuning by developing an information transmission framework. This reimagines the challenge of designing parameter-efficient dense orthogonal matrices as a problem of routing information through a grid-based graph.
    \item Drawing inspiration from butterfly structures in the Cooley-Tukey algorithm, we put forward the Orthogonal Butterfly technique for orthogonal finetuning, which is both parameter-efficient and grounded in the information transmission paradigm.
    \item We offer several theoretical results elucidating why BOFT achieves a marked reduction in the number of trainable parameters without compromising model expressivity or training stability. Furthermore, the matrix factorization characteristic of BOFT inherently supports effective weight interpolation (refer to Figure~\ref{fig:boft_interpolation}).
    \item We are the first to demonstrate the efficacy of orthogonal finetuning in settings that extend beyond controllable text-to-image synthesis~\cite{qiu2023controlling}. Specifically, we adapt BOFT to a suite of tasks spanning both computer vision and natural language processing. In these contexts, BOFT consistently surpasses current leading methods, establishing new standards for parameter-efficiency and model generalization.
\end{itemize}

\section{Related Work}

\textbf{Parameter-efficient finetuning (PEFT)}. The dramatic growth in size and capability of foundation models (\emph{e.g.}, \cite{brown2020language,radford2021learning,rombach2022high,kirillov2023segment}) has led to heightened interest in parameter-efficient strategies for adapting such models to downstream applications~\cite{treviso2023efficient,ding2023parameter,lialin2023scaling}. A wide array of PEFT approaches have been developed~\cite{houlsby2019parameter,aghajanyan2020intrinsic,hulora2022,edalati2022krona,wang2022adamix,gheini2021cross,zaken2022bitfit,guo2020parameter,sung2021training,ansell2022composable,lester2021power,li2021prefix,vu2022spot,he2021towards,mao2021unipelt,karimi2021compacter,liu2022few,sung2022lst,chen2022parameter,jia2022visual,chen2022adaptformer,zhang2022neural,jie2023fact,lian2022scaling,luo2023towards}, but techniques relying on reparameterization~\cite{aghajanyan2020intrinsic,hulora2022,edalati2022krona} are especially pertinent for our research. LoRA~\cite{hulora2022}, for instance, augments the pretrained weight matrices with a low-rank decomposition via the sum of two low-rank matrices, yielding strong performance across natural language processing benchmarks. As LoRA employs a uniform, constant rank for all additional matrices, subsequent studies~\cite{zhang2022adaptive,valipour2022dylora,zhang2023increlora} have explored adaptive rank selection across layers to better utilize the parameter budget. The simplicity of low-rank weight reparameterization has contributed to its widespread adoption~\cite{dettmers2023qlora,zi2023delta,chavan2023one}. Drawing on the role of hyperspherical energy in capturing model generalization~\cite{liu2018learning,liu2021orthogonal}, \cite{qiu2023controlling} introduces orthogonal finetuning (OFT), an alternative mechanism that learns an orthogonal transformation within each layer to adapt diffusion models for text-to-image generation. OFT not only promotes better generalization but also delivers greater training stability than LoRA; however, it usually incurs a higher parameter cost. Consequently, increasing the parameter efficiency of OFT emerges as an important objective. It remains an open question whether OFT can be readily extended to a broader set of adaptation scenarios beyond its initial application to text-to-image diffusion~\cite{qiu2023controlling}. The BOFT method enhances OFT’s parameter economy through butterfly factorization, thereby enabling comprehensive evaluation of orthogonal finetuning across diverse adaptation settings.

\textbf{Butterfly structures}. The radix-2 Cooley-Tukey algorithm efficiently decomposes an $N$-point discrete Fourier transform into two smaller $\frac{N}{2}$-point DFTs in a recursive fashion, yielding a butterfly pattern characterized by a sequence of sparse matrix multiplications (collectively known as a butterfly matrix). Butterfly matrices have served as parameterizations for orthogonal transformations—circumventing the need for pivoting in Gaussian elimination and improving computational efficiency~\cite{parker1995random}—and have been utilized to stabilize recurrent neural network training~\cite{jing2017tunable} and approximate kernels~\cite{munkhoeva2018quadrature}. In addition, \cite{dao2019learning,dao2019kaleidoscope} employ butterfly parameterizations for efficient learning of linear operators. The application of butterfly matrices extends to the scalable training of neural networks~\cite{chen2022pixelated,dao2022monarch}, as well as to fast algorithms for matrix-vector multiplication~\cite{michielssen1996multilevel,o2010algorithm}, data-sparse matrix approximation~\cite{li2015butterfly}, and network information transmission~\cite{klasing1994broadcasting,li2003linear,rajasingh2016transmission}. Diverging from these prior studies, our focus is on leveraging butterfly structures specifically to optimize the parameter usage of OFT.

\section{An Information Transmission View on Orthogonal Finetuning}

We begin by reviewing foundational aspects of OFT. In OFT, the pretrained weight matrix is finetuned by expressing the updated weight matrix as a product between a learnable orthogonal matrix and the fixed pretrained weights. This multiplicative approach contrasts with LoRA, which finetunes weights by adding a low-rank matrix. While LoRA achieves parameter efficiency through low-rank approximations, the standard OFT method adopts a block-diagonal orthogonal parameterization~\cite{qiu2023controlling}, wherein parameter efficiency increases as the size of diagonal blocks decreases. Figure~\ref{fig:comp_lora} illustrates a direct comparison between these two methods. The rationale for leveraging an orthogonal transformation during finetuning lies in its ability to preserve pairwise angular relationships among neurons~\cite{liu2018learning,liu2021orthogonal,qiu2023controlling}, which helps to retain much of the pretrained model's semantic content. More specifically, OFT learns an orthogonal matrix $\bm{R}\in\mathbb{R}^{d\times d}$ for a given pretrained linear layer $\bm{W}^0\in\mathbb{R}^{d\times n}$, updating the forward computation from $\bm{z}=(\bm{W}^0)^\top\bm{x}$ to $\bm{z}=(\bm{R}\bm{W}^0)^\top\bm{x}$, where $\bm{x}\in\mathbb{R}^{d}$ and $\bm{z}\in\mathbb{R}^{n}$ represent input and output vectors, respectively. OFT is zero-initialized by setting $\bm{R}$ to the identity matrix. To maintain the orthogonality of $\bm{R}$ throughout finetuning, the Cayley parameterization is utilized as in~\cite{qiu2023controlling,liu2021orthogonal}; that is, $\bm{R}=(\bm{I}+\bm{Q})(\bm{I}-\bm{Q})^{-1}$ with $\bm{Q}$ being a skew-symmetric matrix satisfying $\bm{Q}=-\bm{Q}^\top$. To further conserve parameters, a block-diagonal scheme is used, representing $\bm{R}$ as $\text{diag}(\bm{R}_1,\bm{R}_2,\cdots,\bm{R}_r)$ where each $\bm{R}_i\in\mathcal{R}^{b\times b}$ for all $i$ and $br=d$. While this block-diagonal structure ensures parameter efficiency, it also imposes an arbitrary assumption that the neuron dimensions (i.e., the columns of $\bm{W}^0$) are split into $r$ groups, with each group’s dimensions transformed independently by different orthogonal matrices. Although empirical results suggest this method is effective, there is no principled reason to assign neuron dimensions to groups solely by index, motivating the exploration of dense orthogonal matrices. Thus, a crucial question emerges: \emph{Is it possible to create a dense orthogonal matrix while retaining parameter efficiency?}

In response to this issue, we suggest constructing a dense orthogonal matrix by multiplying several sparse orthogonal matrices. To create a consistent and accessible framework for analyzing the sparsity structure involved in orthogonal matrix decomposition, we recast the challenge of assembling a dense orthogonal matrix in OFT as analogous to an information transmission task. Concretely, forming a dense matrix $\bm{R}\in\mathbb{R}^{d\times d}$ from the product of $m$ square matrices $\thickmuskip=2mu \medmuskip=2mu \bm{R}=\bm{B}_m\bm{B}_{m-1}\cdots\bm{B}_1$ can be interpreted as the process of passing information through a grid comprising $d\times (m+1)$ nodes, as depicted in Figure~\ref{fig:info_trans}. The underlying rationale for interpreting this process as information flow arises from viewing a $d$-by-$d$ dense square matrix as facilitating dense interactions between two sets of $d$ nodes each. Specifically, the entry $R_{ij}$ of $\bm{R}$ is zero if there is no pathway for information to travel from node $j$ to node $i$, whereas a nonzero $R_{ij}$ implies a possible transmission route. Thus, factoring the dense matrix $\bm{R}$ as a sequence of matrices $\bm{B}_m\bm{B}_{m-1}\cdots\bm{B}_1$ can be equivalently considered as iteratively propagating information over the graphs associated with each $\bm{B}_i$. This information flow starts according to $\bm{B}_1$ and proceeds sequentially, concluding with $\bm{B}_n$. As shown in Figure~\ref{fig:info_trans}, one can consider the specific case of the decomposition $\bm{R}=\bm{B}_5\bm{B}_4\bm{B}_3\bm{B}_2\bm{B}_1$, visualizing both the sparsity patterns and the resulting graph. The depicted graph results from expanding the multiplication steps. In this structure, each $\bm{B}_i$ is interpreted as specifying connections from level $i$ nodes to those in level $i+1$. To be precise, an entry $(j_1,j_2)$ in $\bm{B}_i$ encodes the presence (nonzero) or absence (zero) of a directed edge linking node $j_2$ at the $i$-th level to node $j_1$ at the $(i+1)$-th level. Ensuring that $\bm{B}_5\bm{B}_4\bm{B}_3\bm{B}_2\bm{B}_1$ yields a dense matrix requires that all first-level nodes be capable of transmitting information to each node at the sixth level. If we focus solely on $\bm{R}=\bm{B}_4\bm{B}_3\bm{B}_2\bm{B}_1$, representing flows from initial nodes to those at level five, it becomes apparent, for instance, that information originating at node $1$ does not reach node $3$, signifying that the product has a zero at position $(3,1)$. This relationship between sparsity pattern and connectivity in the associated graphs persists for other decompositions, such as $\bm{R}=\bm{B}_3\bm{B}_2\bm{B}_1$ and $\bm{R}=\bm{B}_2\bm{B}_1$. In general terms, in order for $\bm{R}\in\mathbb{R}^{d\times d}$ to be fully dense, the matrices $\bm{B}_m,\ldots,\bm{B}_1$ must generate a directed edge structure over a $d\times (m+1)$ grid. Each directed edge is constrained to link nodes in consecutive layers (i.e., columns), and the network must permit information from every source node at the first level to reach all destination nodes in the final level.

Adopting the perspective of information transmission provides deeper insights into the sparsity patterns observed in the factorization matrices within OFT. Figure~\ref{fig:blk_diag} illustrates how the original OFT’s matrix $\bm{R}$ exhibits a block-diagonal arrangement. While this structural choice reduces the count of trainable parameters, it inherently prevents $\bm{R}$ from achieving density. The central challenge is thus to assemble a dense orthogonal matrix by leveraging $m$ sparse orthogonal matrices, while minimizing the number of actively trained parameters involved.

From the vantage point of information transmission, two primary requirements emerge: \emph{(i)} \textbf{dense connectivity}, meaning every initial-level node must be reachable by at least one path from every terminal-level node, and \emph{(ii)} \textbf{minimum free edges}, i.e., the overall number of edges must be minimized within the orthogonality constraints. It is important to highlight that orthogonality introduces intricate dependencies across the edges joining adjacent levels. To illustrate, for any matrix $\bm{B}_i$ to maintain full rank (essential for orthogonality), exactly $d$ edges are needed to establish a bijection between nodes in level $i$ and those in level $(i+1)$. This ensures at least $d$ edges link adjacent layers (for instance, see the configuration of four edges in Figure~\ref{fig:info_trans}). These particular $d$ edges are required for orthogonality and should be excluded from the tally of edges, as they are not subject to training (\emph{e.g.}, in a $d\times d$ orthogonal matrix with $d$ non-zero values, only values $\pm 1$ are allowed at those positions).

Owing to the nature of orthogonal matrices, which demand fewer parameters than dense matrices, the orthogonality condition further introduces dependencies among edge presence or absence. For example, in the context of a $2\times 2$ orthogonal matrix, a zero at position $(1,1)$ necessitates a corresponding zero at position $(2,2)$, demonstrating that the omission of one edge may imply the removal of another. Consequently, when considering any feasible configuration of connections, the orthogonality constraint may introduce or eliminate certain edges. Through examination of the non-zero entries in the constructed orthogonal matrix, the information transmission viewpoint becomes especially valuable in OFT, since the focus lies solely on identifying non-zero, trainable parameters in $\bm{R}$—with their specific magnitudes being irrelevant.

A straightforward implementation of a fully connected architecture between two hierarchical levels results in $\mathcal{O}(d^2)$ connections (that is, equivalent to one dense orthogonal matrix), amounting to $d^2-d$ distinct edges (for instance, when $d=4$, this gives 12 edges). In contrast, Figure~\ref{fig:info_trans} demonstrates an example where the matrix is factorized feasibly, requiring only 10 edges overall—fewer than those needed for a single dense orthogonal matrix. This representation allows us to investigate the parameter-efficiency of OFT using graph-theoretic tools, thereby facilitating the generation of efficient factorizations. We draw particular inspiration from a distinctive structure in the Cooley-Tukey algorithm, commonly referred to as butterfly graphs~\cite{cooley1965algorithm}, which are known for connecting $d$ sources to $d$ targets with dense connectivity but utilizing just $\mathcal{O}(d\log d)$ edges. As illustrated in Figure~\ref{fig:info_trans}, the employed topology uses 10 edges for dense connections, whereas the butterfly network topology reduces this to only 8. The following section explores how incorporating the butterfly structure leads to improved parameter efficiency.

\section{Orthogonal Parameterization by Butterfly Factorization}

The butterfly architecture originates from the Cooley-Tukey algorithm, which leverages it to expedite the computation of the fast Fourier transform. In such transformations, an alteration in the frequency space leads to widespread changes in the spatial domain—paralleling our situation, where information at every initial node must be deliverable to each final node. The butterfly configuration has also found significant application as an efficient network topology in computer architecture~\cite{solihin2015fundamentals,li2011linear} to promote effective data exchange. Supposing $k\geq 2$ and $k$ is a power of $2$, we formally define the \emph{butterfly factor} $\bm{B}^F(k)$ as

\begin{equation}\label{eq:but_factor}
\begin{aligned}
    \bm{B}^F(k)=\begin{bmatrix}
    \text{diag}(\bm{d}_1) & \text{diag}(\bm{d}_2) \\
    \text{diag}(\bm{d}_3) & \text{diag}(\bm{d}_4)
    \end{bmatrix}\in\mathbb{R}^{k\times k},
\end{aligned}
\end{equation}

where each vector $\bm{d}_i\in\mathbb{R}^{\frac{k}{2}}$ for $i=1,\ldots,4$. Building on this, let $d=2^N$ for some integer $N$; the corresponding $d$-dimensional \emph{butterfly matrix} $\bm{B}(d)\in\mathbb{R}^{d\times d}$ is recursively defined as

\begin{equation}
\begin{aligned}
    \!\!\bm{B}(d)=\tilde{\bm{B}}(d,d)\!\cdot\!\begin{bmatrix}
    \bm{B}_1(\frac{d}{2})\!\!\!\!\! & \bm{0} \\
    \bm{0}\!\!\!\!\! &\bm{B}_2(\frac{d}{2})
    \end{bmatrix}= \tilde{\bm{B}}(d,d)\tilde{\bm{B}}(d,\frac{d}{2})\cdots\tilde{\bm{B}}(d,2),
\end{aligned}
\end{equation}

Here, $\bm{B}_1(\frac{d}{2})$ and $\bm{B}_2(\frac{d}{2})$ denote two butterfly matrices, each of dimension $\frac{d}{2}$. We introduce the \emph{butterfly component} as $\thickmuskip=2mu \medmuskip=2mu \tilde{\bm{B}}(d,k)=\text{diag}(\bm{B}^F_1(k),\cdots,\bm{B}^F_{d/k}(k))$, forming a block-diagonal matrix of size $d\times d$ with blocks of size $k$. Each $\bm{B}^F_i(k)$, as specified in Equation~\ref{eq:but_factor}, represents a butterfly factor. Utilizing this, we aim to express an orthogonal matrix via the butterfly parameterization. The main requirement is that every multiplicative component $\tilde{\bm{B}}(d,k)$, for all $k$, comprising the butterfly matrix $\bm{B}(d)$, must be orthogonal. To illustrate, consider the block-diagonal structure $\tilde{\bm{B}}(d,2)$ with blocks of size $2$; orthogonality is trivially ensured for each $2\times 2$ block using the Cayley transform (or equivalently a $2$-dimensional rotation), following the approach of \cite{liu2021orthogonal,qiu2023controlling}. Because each butterfly component's non-zero structure is a permutation of that in $\tilde{\bm{B}}(d,2)$, all components can also be straightforwardly constructed as orthogonal matrices. This construction offers an efficient framework to parameterize orthogonal matrices through the composition of numerous $2\times 2$ orthogonal blocks. Building on \cite{chen2022pixelated}, we extend butterfly matrices and define the block butterfly component $\tilde{\bm{B}}^b(d,k)$, where every entry in $\bm{d}_i$ is promoted to a $b\times b$ matrix. To enforce orthogonality for $\tilde{\bm{B}}^b(d,2)$, each $2b\times 2b$ block is parameterized as orthogonal. Since the non-zero pattern for any $\tilde{\bm{B}}^b(d,k)$ with $k>2$ corresponds to a block-wise permutation of that of $\tilde{\bm{B}}^b(d,2)$, these components, too, can be made orthogonal by the same method. When integrated together, the forward computation in BOFT is given by

\begin{equation}
    \begin{aligned}\nonumber
    \bm{z}=\big(\bm{R}(m,b)\cdot\bm{W}^0\big)^\top\bm{x},~~~~\text{s.t.}~\bigg{\{}\bm{R}(m,b)=\prod_{i=1}^m \tilde{\bm{B}}^b_{(d,i)}~~\&~~\underbrace{\big(\tilde{\bm{B}}^b_{(d,j)}\big)^\top\tilde{\bm{B}}^b_{(d,j)}=\tilde{\bm{B}}^b_{(d,j)}\big(\tilde{\bm{B}}^b_{(d,j)}\big)^\top\!=\bm{I}_d}_{\forall j\in [1, m]}\bigg{\}},
    \end{aligned}
\end{equation}

Here, for ease of notation, we write $\tilde{\bm{B}}^b(d,2^{m - i+1})$ as $\tilde{\bm{B}}^b_{(d,i)}$. The symbol $\bm{I}_d$ denotes the $d$-dimensional identity matrix. The orthogonal transformation $\bm{R}(m,b)\in\mathbb{R}^{d\times d}$ is defined as a product of several orthogonal butterfly matrices. To simplify notation, we refer to BOFT with $\bm{R}(m,\frac{b}{2})$ as BOFT($m$,$b$) with $b\geq 2$. In the special case when $\thickmuskip=2mu \medmuskip=2mu m=1$, BOFT($1$,$b$) simplifies to the block-diagonal OFT~\cite{qiu2023controlling} where each block is of size $b$. When $b$ equals $d$, i.e., BOFT($1$,$d$), the method reduces to the original OFT~\cite{qiu2023controlling}, corresponding to using a fully unconstrained orthogonal matrix. If instead we set $m = \log \frac{2d}{b}$ and $b$ as given, BOFT employs the block butterfly matrix $\bm{B}^{\frac{b}{2}}(d)$ for $\bm{R}$, thus forming a dense orthogonal matrix $\bm{R}$. More generally, BOFT($m$,$b$) contains $\frac{1}{2}(b-1)dm$ effective parameters for adapting a linear layer of dimensions $d\times n$. When specifically choosing $m=\log d, b=2$ (the standard butterfly case), BOFT requires $\mathcal{O}(d\log d)$ trainable parameters. On the other hand, standard OFT with a dense orthogonal matrix needs $\mathcal{O}(d^2)$ parameters, whereas the block-diagonal version with $r$ blocks uses $\mathcal{O}(bd)$. Therefore, producing a dense orthogonal matrix with the original OFT demands setting the block size to $b=d$, while BOFT allows for arbitrary $b$ in achieving the same effect.

\textbf{Identity initialization for BOFT}. Commonly in finetuning approaches, initialization closely matches the original pretrained network to ensure minimal divergence at the start of adaptation. LoRA, for example, initializes the additional low-rank matrices as zeros. In the context of BOFT, all butterfly blocks are initialized as identity matrices, meaning that the underlying skew-symmetric matrices in the Cayley transform begin as zero matrices.

\textbf{Multiplicative Dropout for BOFT}. LoRA~\cite{hulora2022} enhances robustness by including a Dropout layer on the low-rank update, a practice that aligns naturally with standard Dropout schemes~\cite{srivastava2014dropout}. However, since BOFT modifies weights multiplicatively, traditional Dropout is not applicable. We address this limitation by designing a multiplicative Dropout tailored for BOFT. The orthogonal structure $\bm{R}(m,b)$ consists of $m$ butterfly matrices, each decomposable into $2b\times 2b$ block-diagonal orthogonal matrices. The proposed multiplicative Dropout works by randomly selecting a portion $p_1$ of the butterfly matrices and, within each, a fraction $p_2$ of their diagonal blocks. These selected blocks are then replaced with identity matrices to implement the dropout effect.

\section{Intriguing Insights and Discussions}

\textbf{Expressivity of BOFT}. The butterfly architecture, together with the use of permutations, is capable of exactly representing several traditional fast linear transforms~\cite{dao2019learning,dao2019kaleidoscope} (\emph{e.g.}, fast Fourier transform, Hadamard transform). However, the extent to which our orthogonal butterfly matrix can approximate an arbitrary orthogonal matrix is not well understood. To investigate this, we carry out a simulation wherein a randomly sampled dense orthogonal matrix~\cite{anderson1987generation} of size $512\times 512$ is approximated. The results depicted in Figure~\ref{fig:para_eff} are averaged over 10 different random seeds. Here, the y-axis measures the approximation error, and the x-axis reflects the number of effective trainable parameters. Curves sharing a color represent BOFT models with identical block sizes; the leftmost mark on each curve represents the approximation error of BOFT($1$,$b$) (that is, the canonical block-diagonal OFT with block size $b$). Overall, BOFT demonstrates higher parameter efficiency than OFT. As an illustration, BOFT($9$,$2$) achieves higher expressivity compared to BOFT($1$,$16$) while utilizing significantly fewer parameters. Generally, BOFT configurations with smaller $b$ values and larger $m$ values tend to be more efficient with respect to parameter count. For instance, BOFT($6$,$4$) matches the approximation performance of BOFT($2$,$16$) but requires far fewer parameters. Thus, the butterfly matrix can be seen as capturing a more structured subset within the orthogonal group, as opposed to the block-diagonal configuration, and consequently, BOFT is provably more expressive than OFT when using the same block size.

\begin{theorem}[Expressivity of BOFT]\label{thm:exp_boft}
    BOFT is strictly more expressive than OFT for an equivalent block size. To allow the butterfly matrix to approximate all orthogonal matrices of dimension $d$, one can form a product of butterfly matrices of the form $\bm{B}_{d-1,1}(d)\bm{B}^\top_{d-1,2}(d)\cdots\bm{B}_{1,1}(d)\bm{B}^\top_{1,2}(d)$, where $\bm{B}_{i,j}(d),\forall i,\forall j$ denote butterfly matrices.
\end{theorem}

Theorem~\ref{thm:exp_boft} highlights a direct pathway for generalizing BOFT—specifically, the terminal orthogonal matrix can be formulated as $\thickmuskip=2mu \medmuskip=2mu \bm{R}^G(m_1,b_1,m_2,b_2,l)=\bm{R}_{l,1}(m_1,b_1)\bm{R}_{l,2}^T(m_2,b_2)\cdots\bm{R}_{1,1}(m_1,b_1)\bm{R}_{1,2}^T(m_2,b_2)$, where $\bm{R}_{i,j}^T(m,b)$ is the orthogonal matrix incorporated in BOFT. Setting $m_1=m_2=\log d$, $b_1=b_2=2$, and $l=d-1$ allows $\bm{R}^G(m_1,b_1,m_2,b_2,l)$ to span the full orthogonal group. This hierarchical matrix construction is also referred to as the kaleidoscope hierarchy~\cite{dao2019kaleidoscope}. It is important to emphasize, however, that increasing model expressivity does not invariably correspond to enhanced finetuning outcomes—full finetuning, despite its theoretical completeness in expressivity, frequently underperforms. Effectively balancing expressivity against regularity is crucial for ensuring robust model generalization during finetuning. BOFT broadens the search space for finetuning parameters via its built-in structural priors, offering a flexible means to negotiate the trade-off between expressivity and regularity.

\textbf{Spectral properties}. Orthogonal finetuning typically outperforms LoRA in maintaining spectral properties, as it exactly retains the spectral norm of the original pretrained weight matrix $\bm{W}^0$. This is evident from the singular value decomposition $\bm{W}^0=\bm{U}\bm{\Sigma}\bm{V}^\top$, where $\bm{U}$ and $\bm{V}$ denote orthogonal matrices and $\bm{\Sigma}$ contains the singular values along its diagonal. Both OFT and BOFT involve multiplying an orthogonal matrix $\bm{R}$ from the left, resulting in updated weights $\bm{R}\bm{U}\bm{\Sigma}\bm{V}^\top$; this operation does not change the maximum singular value—that is, the spectral norm—of $\bm{W}^0$. Prior work demonstrates that maintaining the spectral norm significantly enhances both training stability and the ability to generalize~\cite{miyato2018spectral,yoshida2017spectral}. Further mathematical insights are provided in Appendix~\ref{app:sn_property}.

\textbf{Orthogonal finetuning as learning bilinear similarity}. The BOFT framework can be formulated as learning a bilinear similarity function of the form $\bm{w}^0_i\bm{R}\bm{x}$, with $\bm{w}^0_i$ being the $i$-th neuron (i.e., column vector) from $\bm{W}^0$. In this perspective, BOFT effectively tunes the bilinear similarity matrix $\bm{R}$, which is heavily regularized due to the orthogonality constraint. This establishes a direct relationship with both distance metric learning~\cite{xing2002distance} and classical bilinear forms~\cite{roman2005advanced}.

\textbf{Inductive bias and generalization in BOFT}. Within BOFT, the matrix $\bm{R}(m,b)$ generally resides within a structured subset of the orthogonal group, thus imposing a restriction on the expressivity of the model's hypothesis space. Consequently, BOFT inherently introduces an inductive bias. We posit that the specific bias conferred by butterfly factorization enhances generalization, as it shares a structured motif with a range of traditional linear transforms~\cite{dao2019kaleidoscope}, including the discrete Fourier, discrete sine/cosine, and Hadamard transforms. Additionally, the implementation of sparse matrix factorization within BOFT is likely to introduce further implicit inductive bias~\cite{gunasekar2017implicit,li2018algorithmic,liu2019neural}.

\textbf{Comparison to butterfly-based sparse training}. Numerous prior studies~\cite{dao2019learning,dao2019kaleidoscope,chen2022pixelated,dao2022monarch} have explored sparse neural network training through the butterfly parameterization. These works often involve directly reparameterizing neural network weights as butterfly matrices, proceeding to train models from their initial state. Notably, \cite{dao2022monarch} investigates the finetuning of pretrained weights by initially projecting them onto a modified form of butterfly matrices, followed by optimizing these projected parts for subsequent tasks. In contrast, BOFT introduces an alternative finetuning approach that applies transformations to the model weights using weight matrices shared across layers.

\input{rebuttal-results}

\section{Concluding Remarks and Limitations}

In this work, we present Orthogonal Butterfly, a flexible and parameter-efficient finetuning technique for foundation models, which leverages the butterfly architecture. The main intuition for enhanced parameter efficiency is to express a dense orthogonal matrix as a product of several sparse orthogonal matrices. To streamline the identification of suitable matrix factorizations, we introduce a graphical information transmission paradigm. Within this framework, we observe that butterfly architectures are particularly well-suited for constructing sparse orthogonal matrix factorizations that satisfy our requirements. We empirically demonstrate that BOFT achieves strong performance when applied to the finetuning of large language models, large-scale vision models, and text-to-image generative architectures. Our findings also substantiate the effectiveness of BOFT as a general-purpose finetuning mechanism.

However, BOFT does come with some caveats. Since the ultimate orthogonal matrix in BOFT arises from multiplying several orthogonal matrices together, the training phase incurs a modest increase in computational overhead relative to OFT. Addressing this additional training runtime remains an unresolved challenge. On a positive note, following the finetuning process, the orthogonal matrices learned by BOFT can be consolidated into the pretrained model, leading to no extra latency during inference. Furthermore, it remains an open question whether butterfly networks offer the most optimal structure for information transmission. Our proposed information transmission framework also creates opportunities to incorporate insights from the field of computer networking, where network topology and its impact on efficient information transfer have been thoroughly investigated. This suggests that adopting other network topologies may yield more efficient factorizations for constructing dense orthogonal matrices.